% =====================================================================
% §5 — Implicazioni: il tempo differito mentre il tempo reale è disponibile
%
% Struttura: obiezione nella forma più forte -> Risset (precedente
%            filosofico del ritorno volontario) -> Vaggione (teoria
%            positiva: interazione come triangolo, non come latenza;
%            déclaration d'attribut) -> le tre proposte come conseguenze
%            della postura -> il costo, detto con franchezza
%            (strumenti non neutri) -> chiusa
%
% TODO prima della submission:
%   [x] chiavi BibTeX: FATTO (2026-06-12) — Solomos2003ent->Solomos2003,
%       ArcellaSilvestri2012cim->Arcella2012; le altre erano già reali
%   [ ] verificare pagine delle quote sui PDF: Risset pp. 35, 37;
%       Vaggione 1996 p. 2 e p. 4; entretien pp. 230 e 232;
%       Di Scipio 1995 p. 19; Arcella/Silvestri p. 148
%   [ ] uniformare la resa delle quote (verbatim EN/FR vs parafrasi)
%       allo stile citazionale del resto del paper
%   [ ] il rimando \S\ref{sec:partitura} presuppone il label di §3
% =====================================================================

% ================================================================
\section{Implicazioni: il tempo differito mentre il tempo reale è disponibile}\label{sec:implicazioni}

L'obiezione va formulata nella sua forma più forte: tutto ciò che la
\S\ref{sec:architettura} descrive potrebbe girare in tempo reale su un
laptop, e un granulatore differito ha oggi l'aria dell'anacronismo. La
scelta richiede perciò un argomento, non una giustificazione tecnica ---
e l'argomento ha un precedente esplicito, formulato quando la posizione
era già controcorrente. Per Risset la composizione \emph{«is not --- or
should not be --- a real-time process»}: la notazione riferisce la
musica a una rappresentazione fuori dal tempo, e proprio l'operare non
in tempo reale libera dalla \emph{«arrow of time and its tyranny»},
dai riflessi e dalla fretta; \emph{«writing music implies prediction and
elaboration»}~\cite{Risset1999}. Non è nostalgia del nastro: è la tesi
che un certo tipo di lavoro compositivo --- predizione, elaborazione,
ritorno sul già scritto --- abbia nel differito la propria
configurazione naturale.

Risset da solo, però, esporrebbe il fianco a una lettura del differito
come rinuncia all'interazione. La teoria positiva la fornisce Vaggione.
Nei suoi termini l'interazione non è una proprietà della latenza ma una
relazione a tre --- entrata, uscita, operatore --- e ciò che la
definisce è la presenza dell'operatore nel processo, non la velocità
della risposta~\cite{Solomos2003}. Sul piano della scrittura, ogni
intervento diretto è la \emph{«déclaration d'un attribut»} suscettibile
di essere generalizzata a tutte le istanze successive: scrittura e
algoritmo si imbricano senza che l'una sia infeudata
all'altro~\cite{Vaggione1996}. Letto così, lo YAML del sistema è il
luogo di quell'imbricazione: il valore scritto a mano è la
dichiarazione, l'envelope la sua generalizzazione, la riscrittura dopo
l'ascolto l'azione diretta sul prodotto algoritmico. Il ciclo lungo non
è dunque assenza di interazione ma interazione a un'altra scala
temporale --- e la riduzione dell'interattività a uscita udibile
immediata era già stata rifiutata, in questa stessa sede, trent'anni
fa~\cite{DiScipio1995cim}.

Da questa postura le proposte di \S\ref{sec:tradizione} discendono come
conseguenze, non come dotazioni. Se l'interazione vive fra un ciclo e
l'altro, il ciclo va tenuto corto: di qui la cache, gli stem, il
progetto che si riapre già montato (\S\ref{sec:render}). Se il giudizio
passa dall'osservazione, il risultato dev'essere leggibile: di qui la
partitura (\S\ref{sec:partitura}) --- e l'avvertimento dello stesso
Vaggione contro i tassi di variazione usati come palliativi della visée
figurale è qui assunto alla lettera: i numeri della specifica sono
l'azione, non il giudizio; il giudizio passa dalla figura, osservata e
ascoltata. Se infine la struttura compositiva deve sopravvivere ai
dispositivi --- Risset ancora, contro il rischio di un'arte elettronica
deperibile e senza memoria --- la specifica testuale, esplicita nelle
sue operazioni di base, è precisamente ciò che si conserva e si spedisce.

Ciò che si perde va detto con la stessa franchezza: la performance, il
gesto, lo strumento. Il sistema vi rinuncia deliberatamente --- non è
una voce di roadmap travestita --- e chi cerca lo strumento granulare
troverà altrove, nella linea in tempo reale richiamata in
\S\ref{sec:tradizione}, ciò che qui manca per scelta. Gli strumenti non
sono neutri~\cite{Arcella2012}: questo incorpora la postura
del ciclo lungo, e la incorpora apertamente.

Il tempo differito del 2026 non è quello del 1978: là subìto, qui
scelto. E scelto in quanto abitabile --- perché l'architettura tiene il
ritorno dell'ascolto abbastanza vicino alla scrittura da fare del ciclo
uno spazio di lavoro, non un'attesa.









